\subsection{Atraso de Propagação em Circuitos}

Em circuitos digitais reais, as transições de sinal não ocorrem de forma instantânea. Cada elemento lógico introduz um atraso entre a aplicação de uma entrada e a estabilização da saída, denominado atraso de propagação ($t_p$).

Como consequência, diferentes partes do circuito podem, temporariamente, operar sobre valores inconsistentes de um mesmo sinal. Em circuitos com realimentação, esse fenômeno pode levar a oscilações ou transições espúrias.

O uso de um sinal de clock mitiga esse problema ao restringir as mudanças de estado a instantes discretos, nos quais se assume que os sinais já se estabilizaram.

Considere um circuito combinacional que implementa uma função booleana $f : X \rightarrow Y$. Em um modelo ideal, a saída $y \in Y$ é dada por:

\[
y(t) = f(x(t))
\]

Em um circuito físico, entretanto, a resposta da saída é retardada pelo atraso de propagação. Um modelo temporal mais realista é dado por:

\[
y(t + t_p) = f(x(t))
\]

ou, de forma equivalente,

\[
y(t) = f(x(t - t_p))
\]

Essas expressões indicam que a saída em um dado instante depende de valores passados da entrada, refletindo o comportamento causal do sistema físico.

Esta problemática também está presente nos circuitos de redstone, como será discutido na seção \nameref{sec:dff}.